\documentclass[12pt,a4paper]{mvgl}
\usepackage{epsfig}

\begin{document}

\begin{mvglReport}

% Use \nodate to suppress date or write the day if differs from today's
\date{}

\to{Engin Erzin}

\from{M.~A.~Tu\u{g}tekin~Turan}

\subject{On the Use of MVGL Report Template}

% ---------------------------------------------------------------------------- 
% Main text begins here
% ---------------------------------------------------------------------------- 

This document represents the output from the file "example.tex" once compiled using your favorite \LaTeX compiler.  This file should serve as a good example of the basic structure of a ".tex" file as well as many of the most basic commands needed for typesetting documents involving mathematical symbols and expressions. For more of a description on how each command works, please consult the links found on our course web-page.

\paragraph{Outline}
The remainder of this article is organized as follows.
Section~\ref{pre} gives account of previous work.
Our new and exciting results are described in Section~\ref{res}.
Finally, Section~\ref{conc} gives the conclusions.

\section{Introduction}
\label{Intro}
We use in science. In this report we address the use of dry
erase markers specifically to the science and engineering
performed at the MVGL. Let us begin
with a description of a couple of basic dry erase marker
drawings. We use the circle and the square to make our point.

\subsection{Binomial Theorem}
For any nonnegative integer $n$, we have
$$(1+x)^n = \sum_{i=0}^n {n \choose i} x^i$$

\subsection{Taylor Series}
The Taylor series expansion for the function $e^x$ is given by
\begin{equation}
e^x = 1 + x + \frac{x^2}{2} + \frac{x^3}{6} + \cdots = \sum_{n\geq 0} \frac{x^n}{n!}
\end{equation}

\subsection{Sets}
For any sets $A$, $B$ and $C$, we have
$$ (A\cup B)-(C-A) = A \cup (B-C)$$
\begin{eqnarray*}
(A\cup B)-(C-A) &=& (A\cup B) \cap (C-A)^c\\
&=& (A\cup B) \cap (C \cap A^c)^c \\
&=& (A\cup B) \cap (C^c \cup A) \\
&=& A \cup (B\cap C^c) \\
&=& A \cup (B-C)
\end{eqnarray*}

\subsubsection{Small Circles}
\label{pre}
We also need an example of a sub-subsection to make sure titles and
table of content entries work for those.

And while we are at it, we might as well check paragraph
spacing. Although, we will do that again in Section~\ref{sec:long}.

\begin{enumerate}
\item {\bf First Point (Bold Face)}
\item {\em Second Point (Italic)}
\item {\Large Third Point (Large Font)}
\begin{enumerate}
	\item {\small First Sub-point (Small Font)} 
	\item {\tiny Second Sub-point (Tiny Font)} 
	\item {\Huge Third Sub-point (Huge Font)} 
\end{enumerate}
\item[$\bullet$] {\sf Bullet Point (Sans Serif)}
\item[$\circ$] {\sc Circle Point (Small Caps)} 
\end{enumerate}

\subsection{The Square}
Figure~\ref{fig2} illustrates another famous shape. We will
show some of its uses in later chapters, but we wanted to introduce it
here because this is a good place for another figure.

\begin{figure}[ht]
\centering
\begin{picture}(50,50)(0,0)
\put(0,0){\framebox(50,50){}}
\put(25,25){\circle{1}}
\end{picture}
\caption[The square]{A square is another of the basic
shapes. It is not quite as powerful as the circle. It
has some similarities (note that the four corners all have
the same distance to the center), and has many fine
uses in everyday dry erase marker drawing.}
\label{fig2}
\end{figure}

\subsection{The Dot}
\label{res}
We also need a figure with a short caption. In order to do this,
we introduce the dot. Figure~\ref{fig3} shows the shape of a dot.

\begin{figure}[ht]
\centering
\begin{picture}(50,50)(0,0)
\put(25,25){\circle{5}}
\end{picture}
\caption[The dot]{A simple dot}
\label{fig3}
\end{figure}

\subsection{Tables and Such}
In order to test our class file, we also need to have some tables. One
should be enough for our purposes, so here it is: Table~\ref{tab1}. On
second thought, we need another one to test the list of tables with
multiple entries.

\begin{table}[!bh]
\centering
\caption[Shapes]{This superb table lists a few
of the more important shapes and some of
their properties. Be aware that this condensed list
can by no means describe all the properties or
shapes drawable by dry erase markers.}
\begin{tabular}{|l|c|l|c|}
	\hline\hline
	Name     & Number of & Importance &    Shape    \\
	         &  corners  &            &  \\ \hline
	circle   &     0     & high       & $\bigcirc$  \\
	square   &     4     & medium     & $\diamond$  \\
	triangle &     3     & low        & $\triangle$ \\ \hline
\end{tabular}
\label{tab1}
\end{table}

\section{A Picture}
\begin{center}
\begin{picture}(100,100)(0,0)
\setlength{\unitlength}{1pt}
\put(20,70){\circle{30}}  \put(20,70){\circle*{10}}   % left eye
\put(80,70){\circle{30}}  \put(80,70){\circle*{10}}   % right eye
\put(40,40){\line(1,2){10}} \put(60,40){\line(-1,2){10}} \put(40,40){\line(1,0){20}} % nose
\put(50,20){\oval(80,10)[b]} % mouth
\multiput(0,90)(4,0){10}{\line(1,3){4}}  % left eyebrow
\multiput(100,90)(-4,0){10}{\line(-1,3){4}}  % right eyebrow
\end{picture}
\end{center}

\section{A Long Section}
\label{sec:long}
We need a long chapter to test full-page formatting. Therefore, we
switch to the ancient language of Latin.

\noindent Actually there is one more way, used above; for example, 
{\sc this way}.  The way that you get in and out of environment varies
depending on which kind of environment you want; for example, you use 
\verb|\underline| ``outside'', but \verb|\it| ``inside''; 
notice \underline{this} versus {\it this}.

The real power of \LaTeX\ (for us) is in the math environment. You 
push and pop out of the math environment by typing \verb|$|. For 
example, $2x^3 - 1 = 5$ is typed between dollar signs as
\verb|$2x^3 - 1 = 5$|. Perhaps a more interesting example is $ \sum_{k=1}^N f(t_k) \Delta t. $

You can get a fancier, display-style math 
environment by enclosing your equation with double dollar signs.  
This will center your equation, and display sub- and super-scripts in 
a more readable fashion: $$ \sum_{k=1}^N f(t_k) \Delta t $$

If you don't want your equation to be centered, but you want the nice 
indicies and all that, you can use \verb|\displaystyle| and get your 
formula ``in-line''; using our example this is 
$ \displaystyle \sum_{k=1}^N f(t_k) \Delta t $.  Of 
course this can screw up your line spacing a little bit.

\section{Conclusions}
\label{conc}
Of course, no report would be complete without some conclusions.
This section is where they would go, if we had some.

\begin{table}[htb] 
\centering
\caption{Numbers of Computers on Earth Sciences Network, By Type.}
\begin{tabular}{lr}
	Macintosh                  & 175 \\
	DOS/Windows PC             &  60 \\
	Unix Workstation or server & 110
\end{tabular}
\end{table}

\nocite{*}

% ---------------------------------------------------------------------- %
% References
% ---------------------------------------------------------------------- %
\normalfont 
\bibliographystyle{plain}
\bibliography{exampleBib}

% ---------------------------------------------------------------------- %
% Appendices
% ---------------------------------------------------------------------- %
\clearpage
\renewcommand{\thesection}{\Alph{section}}
\setcounter{section}{0}

\section{Historical Perspective}
Here is a very simple table showing data lined up in columns. Notice that I include the table in a ``center'' environment to display it properly. The title is created simply as another paragraph in the center environment, rather than as part of the table itself.

\subsection{The Past a Long Time Ago}
Simple equations, like $x^y$ or $x_n = \sqrt{a + b}$ can be typeset right
in the text line by enclosing them in a pair of single dollar sign symbols.
Don't forget that if you want a real dollar sign in your text, like \$2000,
you have to use the \verb+\$+ command.

A more complicated equation should be typeset in {\em displayed math\/} mode,
like this:
\[
z \left( 1 \ +\ \sqrt{\omega_{i+1} + \zeta -\frac{x+1}{\Theta +1} y + 1}
\ \right)
\ \ \ =\ \ \ 1
\]
The ``equation'' environment displays your equations, and automatically
numbers them consecutively within your document, like this:
\begin{equation}
\left[
{\bf X} + {\rm a} \ \geq\
\underline{\hat a} \sum_i^N \lim_{x \rightarrow k} \delta C
\right]
\end{equation}

\subsection{The Past More Recently}
Now we have to be a little bit more careful, since records exist from
that time, and some people still alive actually lived back then.

\section{Some Other Appendix}
\begin{enumerate}
\item
The ``enumerate'' environment numbers the list elements, like this.
Items in a list can contain multiple paragraphs.
These paragraphs are appropriately spaced and indented according to their
position in the list.
\begin{itemize}
\item The "itemize" environment sets off list items with "bullets",
like this. Finally, the "description" environment lets you put your own
\begin{description}
\item[X] label on each item, like this ``X''.
\item[If the label is long,] the first line of the item text will
be spaced over to the right as needed.
\end{description} 
\item Of course, lists can be nested, each type up to at least four levels.
One type of list can be nested within another type.
\begin{itemize}
\item Nested lists of the same type will change style of numbering
or ``bullets'' as needed.
\end{itemize}
\end{itemize}
\item Don't forget to close off all list environments with the
appropriate \verb+\end{...}+ command.
Indenting \verb+\begin{...}+, \verb+\item+, and \verb+\end{...}+
commands in the input document according to their nesting level can help
clarify the structure.
\end{enumerate} 

% ---------------------------------------------------------------------- %
% End
% ---------------------------------------------------------------------- %
\end{mvglReport}

\end{document}